% Preámbulo modular para presentaciones Beamer (Modelación Estadística)
% Diseño y paleta institucional del Tecnológico de Monterrey

\documentclass[aspectratio=169,xcolor={dvipsnames,table}]{beamer}

\usepackage[utf8]{inputenc}
\usepackage[spanish,mexico]{babel}
\usepackage{amsmath,amssymb,amsfonts,mathtools}
\usepackage{graphicx}
\usepackage{listings}
\usepackage{booktabs}
\usepackage{tikz}
\usetikzlibrary{matrix,arrows,decorations.pathmorphing,shapes.geometric,positioning}

% Paleta Institucional Tecnológico de Monterrey
\definecolor{TecAzul}{HTML}{0039A6}       % Azul TEC: color primario institucional
\definecolor{TecAzulOscuro}{HTML}{1A2E51} % Azul marino oscuro
\definecolor{TecRojo}{HTML}{EC2661}       % Rojo de acento (paleta miTec)
\definecolor{TecGris}{HTML}{646464}       % Gris neutro para comentarios y subtítulos
\definecolor{TecFondoBloque}{HTML}{F4F6F9}% Fondo suave para bloques

% Configuración de Tema Beamer
\usetheme{Boadilla}
\usecolortheme[named=TecAzulOscuro]{structure}
\setbeamercolor{alerted text}{fg=TecRojo}
\setbeamercolor{palette primary}{bg=TecAzulOscuro,fg=white}
\setbeamercolor{palette secondary}{bg=TecAzul,fg=white}
\setbeamercolor{palette tertiary}{bg=TecAzulOscuro!80!black,fg=white}
\setbeamercolor{title}{fg=TecAzulOscuro,bg=TecFondoBloque}
\setbeamercolor{frametitle}{fg=TecAzulOscuro,bg=TecFondoBloque}

% Bloques personalizados elegantes
\setbeamercolor{block title}{bg=TecAzulOscuro,fg=white}
\setbeamercolor{block body}{bg=TecFondoBloque,fg=black}
\setbeamercolor{block title alerted}{bg=TecRojo,fg=white}
\setbeamercolor{block body alerted}{bg=TecRojo!8,fg=black}
\setbeamercolor{block title example}{bg=TecAzul,fg=white}
\setbeamercolor{block body example}{bg=TecAzul!8,fg=black}

% Redondear bordes y añadir sombra sutil
\setbeamertemplate{blocks}[rounded][shadow=true]
\setbeamertemplate{navigation symbols}{}

% Pie de página limpio con número de diapositiva
\setbeamertemplate{footline}{
  \leavevmode%
  \hbox{%
  \begin{beamercolorbox}[wd=.7\paperwidth,ht=2.25ex,dp=1ex,left]{author in head/foot}%
    \hspace*{2ex}\insertshortauthor~~(\insertshortinstitute) \hfill \insertshorttitle
  \end{beamercolorbox}%
  \begin{beamercolorbox}[wd=.3\paperwidth,ht=2.25ex,dp=1ex,right]{date in head/foot}%
    \insertshortdate{}\hspace*{2em}
    \insertframenumber{} / \inserttotalframenumber\hspace*{2ex}
  \end{beamercolorbox}}%
  \vskip0pt%
}

% Configuración de Listings de Python para visualización pedagógica de código
\lstset{
    language=Python,
    basicstyle=\ttfamily\footnotesize,
    keywordstyle=\color{TecAzulOscuro}\bfseries,
    stringstyle=\color{TecRojo},
    commentstyle=\color{TecGris}\itshape,
    showstringspaces=false,
    breaklines=true,
    frame=lines,
    rulecolor=\color{TecAzulOscuro!30},
    backgroundcolor=\color{TecFondoBloque},
    aboveskip=0.5em,
    belowskip=0.5em,
    literate={á}{{\'a}}1 {é}{{\'e}}1 {í}{{\'i}}1 {ó}{{\'o}}1 {ú}{{\'u}}1
             {Á}{{\'A}}1 {É}{{\'E}}1 {Í}{{\'I}}1 {Ó}{{\'O}}1 {Ú}{{\'U}}1
             {ñ}{{\~n}}1 {Ñ}{{\~N}}1 {¿}{{?`}}1 {¡}{{!`}}1
}

% Metadatos institucionales por defecto
\title[Modelación Estadística]{Modelación Estadística para Ciencia de Datos e Ingeniería}
\author[J. Castillo Colmenares]{Juliho David Castillo Colmenares}
\institute[Tec de Monterrey]{Tecnológico de Monterrey \\ Escuela de Ingeniería y Ciencias}
\date{\today}

% Comandos y macros estadísticas compartidas para las presentaciones Beamer (Metropolis)

% Conjuntos numéricos básicos
\newcommand{\R}{\mathbb{R}}
\newcommand{\N}{\mathbb{N}}
\newcommand{\Q}{\mathbb{Q}}
\newcommand{\Z}{\mathbb{Z}}
\newcommand{\set}[1]{\left\{ #1 \right\}}
\newcommand{\abs}[1]{\left|#1\right|}
\newcommand{\norm}[1]{\left\| #1 \right\|}

% Comandos estadísticos y probabilísticos del curso
\newcommand{\Var}{\operatorname{Var}}
\newcommand{\cov}{\operatorname{Cov}}
\newcommand{\corr}{\rho}
\newcommand{\comb}[2]{\begin{pmatrix} #1 \\ #2 \end{pmatrix}}
\newcommand{\s}{\sigma}
\newcommand{\particion}{\mathfrak{P}}
\newcommand{\card}[1]{\#\left( #1 \right)}
\newcommand{\seq}[3]{\set{#1}_{#2}^{#3}}
\newcommand{\E}{\mathbb{E}}
\newcommand{\Prob}{\mathbb{P}}

% Entornos pedagógicos y cajas en español académico natural
\newenvironment{discusion}{%
  \begin{alertblock}{Pregunta de Discusión}%
}{%
  \end{alertblock}%
}

\newenvironment{reflexion}{%
  \begin{alertblock}{Reflexión para el Aula}%
}{%
  \end{alertblock}%
}

\newenvironment{paradoja}{%
  \begin{alertblock}{Paradoja de Exploración}%
}{%
  \end{alertblock}%
}

\newenvironment{motivacion}{%
  \begin{exampleblock}{Motivación: Ciencia de Datos en la Práctica}%
}{%
  \end{exampleblock}%
}

\newenvironment{retoclase}{%
  \begin{block}{Reto Rápido en Equipo}%
}{%
  \end{block}%
}


\title[Notación de Conjuntos]{Sección 2.2: Notación de Conjuntos y Particiones}
\subtitle{El Lenguaje Algebraico y Causal del Azar}
\author{Julio César Hernández Castro}
\institute{Tecnológico de Monterrey}
\date{}

\begin{document}

% Slide 1: Portada
\begin{frame}[plain]
  \titlepage
\end{frame}

% Slide 2: Hoja de Ruta e Indicador de Posición
\begin{frame}{Hoja de Ruta — Capítulo 02: Teoría de Probabilidad}
  ¿Dónde estamos en nuestro recorrido por la teoría de probabilidad? \pause
  \begin{itemize}
    \item \textbf{02.01 Introducción a la probabilidad} \emph{(Completado)} \pause
    \item \color{TecRojo}\textbf{02.02 Notación de conjuntos y particiones \emph{(Hoy)}}\color{black} \pause
    \item \textbf{02.03 Fundamentos y axiomas de probabilidad} \emph{(Siguiente sesión)} \pause
    \item \textbf{02.04 Probabilidad condicional e independencia} \pause
    \item \textbf{02.05 Teorema de Bayes} \pause
    \item \textbf{02.06 Muestreo aleatorio}
  \end{itemize}
  \vspace{0.6em}
  \begin{block}{Objetivo de la Sección 2.2}
    Estandarizar las operaciones de conjuntos y utilizar \emph{particiones} como rompecabezas algebraicos para cuantificar elementos y evitar el conteo duplicado en análisis causales de datos.
  \end{block}
\end{frame}

% Slide 3: Motivación / Dilema de las Encuestas
\begin{frame}{El Dilema del Conteo en Ciencia de Datos}
  \begin{motivacion}
    En una encuesta a $100$ usuarios de una app bancaria, $60$ utilizan dispositivos \textbf{iOS}, $45$ utilizan \textbf{Android} y $30$ utilizan ambos plataformas para sus transacciones. Si un directivo suma $60 + 45 = 105$, supera el número total de encuestados. ¿Por qué ocurre esta paradoja de conteo?
  \end{motivacion}
  
  \vspace{0.4em}
  \pause
  \begin{itemize}
    \item El error de \textbf{conteo duplicado} (\emph{double counting}) surge al ignorar la estructura algebraica de las intersecciones en conjuntos que no son mutuamente excluyentes. \pause
    \item Para medir probabilidades con rigor exacto, necesitamos descomponer el espacio muestral en un \textbf{rompecabezas de particiones disjuntas}.
  \end{itemize}
\end{frame}

% Slide 4: Álgebra de Conjuntos y Operaciones Básicas
\begin{frame}{Operaciones Fundamentales con Conjuntos}
  Sea un universo o espacio muestral $S$ y subconjuntos $A, B \subset S$:
  
  \vspace{0.2em}
  \begin{itemize}
    \item \textbf{Unión ($A \cup B$):} Elementos en $A$, en $B$, o en ambos: $\{x \in S \mid x \in A \text{ o } x \in B\}$. \pause
    \item \textbf{Intersección ($A \cap B$):} Elementos simultáneamente en ambos: $\{x \in S \mid x \in A \text{ y } x \in B\}$. \pause
    \item \textbf{Diferencia ($A \setminus B$):} Elementos en $A$ que no pertenecen a $B$: $\{x \in S \mid x \in A \text{ y } x \notin B\}$. \pause
    \item \textbf{Complemento ($A'$):} Elementos de $S$ fuera de $A$: $\{x \in S \mid x \notin A\} = S \setminus A$. \pause
    \item \textbf{Diferencia Simétrica ($A \triangle B$):} Elementos en $A$ o $B$, pero no en ambos: $(A \setminus B) \cup (B \setminus A)$.
  \end{itemize}
\end{frame}

% Slide 5: El Concepto de Partición (Rompecabezas)
\begin{frame}{¿Qué es una Partición de un Conjunto?}
  Una \textbf{partición} de un conjunto $S$ es una colección finita de subconjuntos $\{P_i\}_{i=0}^{N}$ que divide el espacio total en piezas exactas, como un rompecabezas:
  
  \vspace{0.4em}
  \begin{block}{Definición Formal de Partición}
    La colección $\{P_i\}_{i=0}^{N}$ es una partición de $S$ si y solo si cumple dos condiciones:
    \begin{enumerate}
      \item \textbf{Mutuamente Disjuntos:} $P_i \cap P_j = \emptyset$ siempre que $i \neq j$. Ninguna pieza del rompecabezas se traslapa con otra. \pause
      \item \textbf{Cobertura Total:} $\bigcup_{i=0}^{N} P_i = S$. Al unir todas las piezas se reconstruye el espacio total sin dejar vacíos.
    \end{enumerate}
  \end{block}
\end{frame}

% Slide 6: Partición Estándar de 2 Conjuntos (4 Regiones)
\begin{frame}{Partición Estándar para Dos Conjuntos}
  Para dos eventos $A, B \subset S$, cada elemento puede pertenecer o no a cada uno de ellos ($2^2 = 4$ regiones elementales disjuntas):
  
  \vspace{0.3em}
  \begin{itemize}
    \item $E_0 = A' \cap B'$ \quad (Elementos en ninguna de las dos categorías). \pause
    \item $E_1 = A' \cap B$ \quad (Elementos exclusivos del conjunto $B$). \pause
    \item $E_2 = A \cap B'$ \quad (Elementos exclusivos del conjunto $A$). \pause
    \item $E_3 = A \cap B$ \quad (Elementos en la intersección simultánea).
  \end{itemize}
  
  \vspace{0.3em}
  \pause
  \begin{block}{Regla Aditiva Exacta sobre Particiones}
    Dado que $E_i \cap E_j = \emptyset$, el conteo cardinal es estrictamente lineal:
    \[ \card{S} = x_0 + x_1 + x_2 + x_3, \quad \text{donde } x_k = \card{E_k} \]
  \end{block}
\end{frame}

% Slide 7: Conexión con Sistemas de Ecuaciones Lineales
\begin{frame}{Resolución Algebraica mediante Sistemas Lineales}
  Cualquier problema de conteo sobre $A, B$ se traduce en un sistema lineal sobre las incógnitas de la partición ($x_0, x_1, x_2, x_3$):
  
  \vspace{0.3em}
  \begin{center}
    \begin{tabular}{l l l}
      \textbf{Conjunto o Condición} & \textbf{Suma sobre Partición} & \textbf{Ejemplo iOS/Android} \\
      \hline
      Espacio Muestral ($S$) & $x_0 + x_1 + x_2 + x_3$ & $= 100$ \\
      Conjunto $A$ (\emph{iOS}) & $x_2 + x_3$ & $= 60$ \\
      Conjunto $B$ (\emph{Android}) & $x_1 + x_3$ & $= 45$ \\
      Intersección ($A \cap B$) & $x_3$ & $= 30$ \\
      \hline
    \end{tabular}
  \end{center}
  
  \vspace{0.3em}
  \pause
  Al sustituir $x_3 = 30$, obtenemos de inmediato $x_2 = 30$, $x_1 = 15$ y $x_0 = 25$. ¡La suma exacta es $30 + 30 + 15 + 25 = 100$!
\end{frame}

% Slide 8: Partición de 3 Conjuntos (8 Regiones)
\begin{frame}{Partición Estándar para Tres Conjuntos}
  Al introducir un tercer conjunto $C \subset S$, el espacio se divide en $2^3 = 8$ piezas elementales que podemos enumerar en código binario:
  
  \vspace{0.2em}
  \begin{columns}[T]
    \begin{column}{0.48\textwidth}
      \small
      \begin{itemize}
        \item $E_0 = A' \cap B' \cap C'$ \quad ($000$)
        \item $E_1 = A' \cap B' \cap C$ \quad ($001$)
        \item $E_2 = A' \cap B \cap C'$ \quad ($010$)
        \item $E_3 = A' \cap B \cap C$ \quad ($011$)
      \end{itemize}
    \end{column}
    \begin{column}{0.48\textwidth}
      \small
      \begin{itemize}
        \item $E_4 = A \cap B' \cap C'$ \quad ($100$)
        \item $E_5 = A \cap B' \cap C$ \quad ($101$)
        \item $E_6 = A \cap B \cap C'$ \quad ($110$)
        \item $E_7 = A \cap B \cap C$ \quad ($111$)
      \end{itemize}
    \end{column}
  \end{columns}
  
  \vspace{0.3em}
  \pause
  \begin{center}
    \small Esto fundamenta la famosa \textbf{Fórmula de Inclusión-Exclusión}:
    \[ \card{A \cup B \cup C} = \sum \card{A_i} - \sum \card{A_i \cap A_j} + \card{A \cap B \cap C} \]
  \end{center}
\end{frame}

% Slide 9A: Laboratorio Python (Parte 1: Operaciones Básicas)
\begin{frame}[fragile]{Laboratorio en Python: Operaciones Básicas}
  El script \texttt{02.02\_conjuntos\_particiones.py} modela álgebra de eventos usando \texttt{set}:
  \vspace{0.2em}
  {\lstset{basicstyle=\fontsize{6.5pt}{7.5pt}\selectfont\ttfamily}
  \lstinputlisting[firstline=1, lastline=16]{../../code/02_teoria_probabilidad/02.02_conjuntos_particiones.py}}
\end{frame}

% Slide 9B-1: Laboratorio Python (Parte 2: Planteamiento del Sistema)
\begin{frame}[fragile]{Laboratorio en Python: Ecuaciones de Partición}
  Planteamiento de las 4 regiones elementales como incógnitas del sistema:
  \vspace{0.1em}
  {\lstset{basicstyle=\fontsize{6.5pt}{7.5pt}\selectfont\ttfamily}
  \lstinputlisting[firstline=18, lastline=35]{../../code/02_teoria_probabilidad/02.02_conjuntos_particiones.py}}
\end{frame}

% Slide 9B-2: Laboratorio Python (Parte 3: Matriz Numpy)
\begin{frame}[fragile]{Laboratorio en Python: Matriz del Sistema Lineal}
  Definición matricial $M x = b$ usando \texttt{numpy.array}:
  \vspace{0.1em}
  {\lstset{basicstyle=\fontsize{6.5pt}{7.5pt}\selectfont\ttfamily}
  \lstinputlisting[firstline=37, lastline=46]{../../code/02_teoria_probabilidad/02.02_conjuntos_particiones.py}}
\end{frame}

% Slide 9C: Laboratorio Python (Parte 4: Solución y Verificación)
\begin{frame}[fragile]{Laboratorio en Python: Solución e Inspección}
  Resolución matricial con \texttt{numpy.linalg.solve} y validación de la suma cardinal:
  \vspace{0.1em}
  {\lstset{basicstyle=\fontsize{6.5pt}{7.5pt}\selectfont\ttfamily}
  \lstinputlisting[firstline=47, lastline=54]{../../code/02_teoria_probabilidad/02.02_conjuntos_particiones.py}}
\end{frame}

% Slide 9D: Laboratorio Python (Salida en Consola)
\begin{frame}[fragile]{Laboratorio en Python: Salida en Terminal}
  Consola de resultados verificando las cardinalidades exactas de la partición:
  \vspace{0.1em}
  \begin{block}{Salida Estándar en Terminal}
    \fontsize{6.5pt}{7.5pt}\selectfont
    \begin{verbatim}
--- 1. Operaciones Básicas de Conjuntos ---
Espacio Muestral (S):  [1, 2, 3, 4, 5, 6, 7, 8, 9, 10]
Evento A (Pares):      [2, 4, 6, 8, 10]
Evento B (<= 5):       [1, 2, 3, 4, 5]
Unión (A union B):     [1, 2, 3, 4, 5, 6, 8, 10]
Intersección (A & B):  [2, 4]
Diferencia (A \ B):    [6, 8, 10]
Complemento (A'):      [1, 3, 5, 7, 9]

--- 2. Resolución Algebraica vía Partición Estándar (4 Piezas) ---
x3 (|I & F|   - Ambos idiomas):           30 personas
x2 (|I \ F|   - Solo Inglés):            30 personas
x1 (|F \ I|   - Solo Francés):           15 personas
x0 (|I' & F'| - Ninguno de los dos):      25 personas
Total verificado de la partición:        100 personas
    \end{verbatim}
  \end{block}
\end{frame}

% Slide 10: Ejercicios del Libro (Tiered 3-3-2-2) - Fundamental
\begin{frame}{Ejercicios del Libro — Nivel 1: Fundamental}
  \begin{block}{Problema 2.1.2 (Congreso Internacional)}
    De $100$ asistentes a un congreso, $60$ hablan inglés ($I$), $45$ hablan francés ($F$) y $30$ hablan ambos idiomas ($I \cap F$). Determina:
    \begin{enumerate}
      \item $\card{I \cup F}$, el total que habla al menos uno de los dos idiomas. \pause
      \item $\card{I \setminus F}$, los que hablan estrictamente solo inglés. \pause
      \item $\card{(I \cup F)'}$, el número de asistentes que no habla ni inglés ni francés.
    \end{enumerate}
  \end{block}
  
  \vspace{0.3em}
  \pause
  \begin{alertblock}{Solución Inmediata vía Partición}
    Con $x_3=30$, $x_2=30$, $x_1=15$ y $x_0=25$:
    1) $\card{I \cup F} = 75$; \quad 2) $\card{I \setminus F} = 30$; \quad 3) $\card{(I \cup F)'} = 25$.
  \end{alertblock}
\end{frame}

% Slide 11: Ejercicios del Libro (Tiered 3-3-2-2) - Operativo
\begin{frame}{Ejercicios del Libro — Nivel 2: Operativo}
  \begin{block}{Problema 2.1.4 (Extensión a Tres Subconjuntos)}
    Extiende la construcción de una partición al caso de tres subconjuntos $A, B, C \subset S$. ¿Cuántos elementos tiene la partición $\particion(A,B,C)$ y qué propiedad algebraica garantiza que no haya traslapos?
  \end{block}
  
  \vspace{0.4em}
  \pause
  \begin{itemize}
    \item \textbf{Número de Elementos:} Cada conjunto duplica las regiones anteriores ($2^3 = 8$ regiones elementales disjuntas $E_0, E_1, \dots, E_7$). \pause
    \item \textbf{Garantía de Disjunción:} Para cualquier $i \neq j$, sus representaciones binarias difieren en al menos un bit (ej. estar en $A$ vs. estar en $A'$), cuya intersección es por definición el conjunto vacío ($\emptyset$).
  \end{itemize}
\end{frame}

% Slide 12: Ejercicios del Libro (Tiered 3-3-2-2) - Analítico
\begin{frame}{Ejercicios del Libro — Nivel 3: Analítico}
  \small
  \begin{block}{Problema 2.1.8 (Inclusión-Exclusión para 3 Conjuntos)}
    Demuestra la identidad para la unión de tres conjuntos:
    \begin{align*}
      \card{A \cup B \cup C} = \card{A} + \card{B} + \card{C} - \card{A \cap B} - \card{A \cap C} \\
      - \card{B \cap C} + \card{A \cap B \cap C}
    \end{align*}
  \end{block}
  
  \pause
  \textbf{Demostración vía Particiones ($x_k = \card{E_k}$):}
  \begin{itemize}
    \item $\card{A} + \card{B} + \card{C}$ cuenta intersecciones dobles $2$ veces y la triple $3$ veces. \pause
    \item Restar las dobles compensa los traslapos pero elimina $x_7$ ($3 - 3 = 0$). \pause
    \item Por ello, sumamos $+ \card{A \cap B \cap C} = +x_7$ para obtener la cobertura exacta.
  \end{itemize}
\end{frame}

% Slide 13: Ejercicios del Libro (Tiered 3-3-2-2) - Desafiante
\begin{frame}{Ejercicios del Libro — Nivel 4: Desafiante}
  \small
  \begin{block}{Problema 2.1.10 (Absurdos Censales en Epidemiología)}
    En $500$ pacientes, un informe dice: $\card{S_1}=300, \card{S_2}=250, \card{S_3}=200$; $\card{S_1 \cap S_2}=180, \card{S_1 \cap S_3}=150, \card{S_2 \cap S_3}=130$; $\card{S_1 \cap S_2 \cap S_3}=30$. ¿Por qué es falso?
  \end{block}
  
  \pause
  \begin{alertblock}{Análisis de la Región de Complemento ($x_0$)}
    Intersecciones binarias simples: $\card{S_1 \cap S_2 \cap S_3'} = 150$, $\card{S_1 \cap S_3 \cap S_2'} = 120$, $\card{S_2 \cap S_3 \cap S_1'} = 100$.
    \pause
    Suma elemental: $150+120+100+30 + \text{simples} > 500$. Por tanto $x_0 < 0$ (imposible).
  \end{alertblock}
\end{frame}

% Slide 14: Síntesis y Próximo Paso
\begin{frame}{Resumen de la Sección 2.2}
  \begin{block}{Conclusiones Clave}
    \begin{itemize}
      \item Toda operación sobre conjuntos no excluyentes debe transformarse en una \textbf{partición disjunta} para garantizar sumas aditivas sin sesgos. \pause
      \item Los sistemas de ecuaciones lineales sobre particiones ($Mx = b$) son la herramienta computacional óptima para validar encuestas y censos.
    \end{itemize}
  \end{block}
  
  \vspace{0.6em}
  \pause
  {\large \color{TecAzulOscuro}\textbf{Siguiente Sesión: 02.03 Fundamentos y Axiomas de Probabilidad}}
  \vspace{0.3em}
  \begin{itemize}
    \item Con la notación de conjuntos dominada, introduciremos la función de medida probabilística $P(A)$ y los tres axiomas inmortales de Kolmogorov.
  \end{itemize}
\end{frame}

\end{document}
